% New independent local approximation theorem; no numerical integral is certified here.
\subsection{Joint finite approximation of a causal lower-tree policy}
\label{sec:lower-start-joint-finite-approximation}

This is an approximation result for the already specified signed
lower-tree construction. Its arithmetic premise is the finite-policy
compiler, with its fixed-cutoff, fixed-margin signed-atom distribution
estimate and ordinary residual-density equalities. It is not a proof
of a floating integral, and does not strengthen that arithmetic premise.

Put $\nu=249/448$, $0<w\le h=19/125$, and $\Lambda=153/256$.
An ideal policy consists of causal caps
$C_r(t_1,\ldots,t_r)$, $r=3,4,5$, with
\[
 \nu\le C_3\le C_4\le C_5\le\Lambda.
\]
Each cap is either inherited from its parent or the endpoint of one
of finitely many licensed branches. Branch endpoints and gates are
piecewise affine on a finite polyhedral partition; the present
$v$, $\Phi$ and selected-subproduct branches have this form.
The ideal policy itself, including its branch labels, is required
to be piecewise affine/constant on a finite polyhedral refinement,
as for the finite maxima with deterministic tie rules considered
here. Arbitrary measurable switching between candidate labels is
outside this statement.
An active new branch must have a nonempty strict interval between
all its gates, including the parent cap, and its endpoint.
The branch label and its proof of support are part of the state,
not merely the scalar cap.

The ideal admitted domains are
\begin{align*}
 E_3&=\{0<t_3<t_2<t_1<w:\ t_1+3t_2<\nu\},\\
 E_4&=\{(t_1,\ldots,t_4):(t_1,t_2,t_3)\in E_3,
       \ 0<t_4<t_3,\ S_3+3t_4<C_3\},\\
 E_5&=\{(t_1,\ldots,t_5):(t_1,\ldots,t_4)\in E_4,
                                      \ 0<t_5<t_4\}.
\end{align*}
Here $S_r=\sum_{j\le r}t_j$. The fifth square is automatic:
$S_4+2t_5<S_3+3t_4<C_3\le C_4$. This also holds for a finite
policy when its fourth historical cube is checked at its actual
frozen parent cap.

Write $d\omega_r=dt_1\cdots dt_r/(t_1\cdots t_{r-1}t_r^2)$.
The ideal signed-density coefficient is its ordinary starting
coefficient plus the three exact telescoping integrals
\begin{align*}
 I_3&=4\int_{E_3}
  [\bar F((\nu-S_3)/t_3)-\bar F((C_3-S_3)/t_3)]\,d\omega_3,\\
 I_4&=4\int_{E_4}
  [\bar f((C_4-S_4)/t_4)-\bar f((C_3-S_4)/t_4)]\,d\omega_4,\\
 I_5&=4\int_{E_5}
  [\bar F((C_4-S_5)/t_5)-\bar F((C_5-S_5)/t_5)]\,d\omega_5.
\end{align*}
All three integrands are nonnegative. They belong to one causal
policy; a fifth integral with ideal $C_4$ cannot instead be appended
to an unrelated finite fourth certificate.

\paragraph{Finite-approximation statement.}
For every $\eta>0$ there is one finite nested prefix partition and
one legal buffered policy $\Pi$, chosen independently of $N$, whose
ordinary raw-density coefficient differs from the ideal coefficient
by less than $\eta$. Consequently the finite-policy compiler gives
a lower bound with coefficient at least the ideal coefficient minus
$2\eta$, for every sufficiently large even $N$.
The second $\eta$ is reserved for compilation and arithmetic errors.
If the old endpoint $\nu$ itself needs distribution slack, use the
ordinary baseline $\nu_\rho=\nu-\rho$ in the finite construction and
let $\rho\downarrow0$ as part of this approximation. No exact endpoint
distribution assertion is being silently added.

\paragraph{Proof on a compact truncated domain.}
First fix $t_r\ge\epsilon>0$ for every active rank. The ordinary
functions are continuous and bounded on all attained argument
intervals. Old admission gives the third and fifth upper arguments
greater than one and the fourth lower arguments greater than two.
Every argument is at most $\Lambda/\epsilon$. Thus each integrand
is bounded on this compact domain and is uniformly continuous
in the cap variables away from the finite policy walls.

Refine all affine pieces and strict branch tests to a finite
polyhedral arrangement. For two genuinely different affine expressions,
their tie is a hyperplane and has measure zero. Identically equal
expressions are merged or assigned a fixed label; they are not
declared a boundary. In particular the positive-measure region
$C_4=C_3$ is retained with an explicit \emph{inherit} label.
It can support a positive fifth improvement and must not be discarded.
The same convention applies when $C_3=\nu$ or $C_5=C_4$.

Trim neighborhoods of the genuine walls, old cube boundaries,
ordering boundaries and nonidentical gate equalities until their
total weighted integral is small. A gate endpoint identical to its
lower gate on an open chamber is not an available new action:
there is no strictly buffered level between them. This is distinct
from the legitimate parent-inheritance label.

Now partition the first three coordinates into sufficiently small
rational boxes. A finite third state must satisfy its whole-box old
history $t_{1,+}+3t_{2,+}<\nu_\rho$ (and its first square), rather
than retaining an unbuffered endpoint history. On a retained active
chamber choose the licensed
branch and freeze its whole-box infimum minus a small positive
buffer. Keep $\nu_\rho$ on inherited or unused cells. For each such
fixed parent, partition its product with the fourth-coordinate interval;
never recompute its parent level using the fourth coordinate.
Recheck $S_{3,+}+3t_{4,+}<\lambda_3$ with that actual parent value.
On an inherited fourth chamber set $\lambda_4=\lambda_3$ exactly;
on an active chamber use the selected branch infimum minus a buffer,
and recheck every full and selected gate. Finally partition the fifth
coordinate under each existing fourth state, including the states
with no fourth improvement, and apply the same procedure there.
Later refinements of a parent box retain its already frozen cap.

This construction is causal and finite. It is not an arbitrary
five-dimensional partition whose nominal third level can vary with
$t_4$ or $t_5$. Nor does it delete the projection onto parent space
of every child boundary; such a projection can have full measure.
Child boundary cells are handled at their own rank, while their
parents remain available with ordinary descendants.

On the compact trimmed active chambers all relevant strict margins
have a positive minimum. The affine endpoints are uniformly
continuous. Hence sufficiently small boxes and buffers preserve
every gate and make the frozen caps uniformly close to $C_r$.
Equivalently, along a sequence of these constructions the caps
converge almost everywhere and remain below the corresponding ideal
caps. At an admissible fourth point the positive margin
$C_3-S_3-3t_4$ ensures eventual admission under the approximating
parent cap. Its zero set is a graph in $t_4$, since $C_3$ is independent
of $t_4$, and therefore has measure zero. The fifth historical square
follows from the preserved fourth history and the ordering.
The domain indicators and all three kernels consequently converge
almost everywhere. Bounded convergence proves convergence of their
three compact integrals and of the ordinary baseline coefficient.

It is neither necessary nor generally correct to claim that the
individual fourth and fifth integral approximations are monotone:
lowering a parent cap changes both a kernel and its admissible domain.
The proof is convergence of exact telescoping differences, not
pointwise monotonicity of signed divisor weights.

\paragraph{Removing the positive cutoff.}
Put $d_0=\nu-3h>0$. The ordinary estimates
$\bar F(s)-e^{-\gamma}/2=O(s^{-2})$ and
$e^{-\gamma}/2-\bar f(s)=O(s^{-2})$ suffice. At rank three,
$\nu-S_3\ge d_0$. At rank four the historical cube gives
\[
 C_3-S_4>\tfrac23(C_3-S_3)\ge\tfrac23d_0.
\]
The same lower bound holds for $C_4-S_4$ before the fifth prime
is removed. Thus, for sufficiently small last coordinate, each
function difference cancels its $t_r^{-2}$ kernel singularity.
Integration over the remaining ordered coordinates gives tails
\[
 I_r(t_r<\epsilon)
   =O\bigl(\epsilon[1+\log^{r-1}(h/\epsilon)]\bigr),
 \qquad r=3,4,5.
\]
The bounds are uniform for small baseline buffers, say
$\rho<d_0/2$, after replacing $d_0$ by $d_0/2$.
This proves absolute convergence and permits choosing the cutoff
before the finite mesh. It is not an exchange of the arithmetic
$N$ limit with an unbounded family of strategies.

\paragraph{Arithmetic quantifiers and what remains to be certified.}
The compiler premise needed here is precise: for each fixed finite
buffered policy, each fixed positive rough cutoff and fixed choice
of the auxiliary small-prime parameters, its signed convolution
atoms satisfy the distribution estimate with constants allowed to
depend on that policy, uniformly in the moving residue associated
with every even $N$. The fixed-label factorization and beta/sign
compilation are the ones in the existing joint audits, with the
lower parity choices proved in the starting-lower note. There are
finitely many atoms; their error constants can therefore be summed
and their thresholds maximized. No bound uniform as the number of
boxes tends to infinity is needed or asserted.

The order is: choose an error budget, a positive cutoff, a finite
partition and cap buffers, then the fine-box/beta/HB and local-scale
parameters compatible with their smallest margin; only then take
$N$ sufficiently large. For a sequence of policies $\Pi_m$, one
must use $N\ge N_0(\Pi_m)$ for the chosen fixed $m$. An estimate
with $\Pi=\Pi(N)$ has not been proved.

Accordingly, a \emph{rigorous} lower bound $c_*$ for the complete
ideal functional, strictly above a target, would suffice: choose
$\eta$ below the gap, fix the resulting finite policy, and apply
the preceding arithmetic quantifiers. It would not be logically
necessary to store millions of explicit policy boxes. But the
integral lower bound, the inherited distribution theorem and any
other terms of that complete functional still need proofs.
Sobol samples, floating tables and agreement of seeds do not supply
that missing numerical integral certificate. This lemma alone
therefore changes no presently certified global constant.

\paragraph{Concrete failure avoided by the construction.}
If a finite parent has $\lambda_3<C_3$, then, whenever it lies below
$t_3$, the choice
$t_4=(C_3+\lambda_3-2S_3)/6$ satisfies
$S_3+3t_4=(C_3+\lambda_3)/2>\lambda_3$ while passing the ideal
fourth history. Using it in that finite parent's fourth gain is
invalid. The discrepancy is a thin region that vanishes in the
joint approximation; it cannot be counted before the histories are
retested. This is why separate ideal-layer integrals cannot be
combined with an arbitrary already frozen ancestor certificate.
